\documentclass[../main.tex]{subfiles}
\begin{document}
\chapter{Minimising with Constraints}
\section{Motivation}
Suppose we have a surface $S \subseteq \R^3$ and a function $f: \R^3 \to \R$.
How do we determine the minimum and maximum value of $f$ \textbf{on $S$ specifically}, not just over all of $\R^{3}$.

For example, if $S$ represents the surface of a hill in 3D space and $f$ the temperature at a given point, then we might want to find the hottest and coldest points that lie on the hill and not just across the entire domain of $f$.

In this case, we cannot just solve $\nabla f = \vec{0}$ as the extrema could be on the boundary of $S$ and so would not have $\nabla f = \vec{0}$.
\section{Lagrange Multipliers}
Let $(u, v)$ be a parameterisation of $S$, i.e. $\vec{x} = \vec{x}(u, v)$ on $S$.
We then wish to extremise $f(\vec{x}(u, v))$.
Since this can now be thought of a function of $(u, v)$, to find the stationary points we just equate both partial derivatives with 0:
\begin{align*}
  0 &= \pderiv{}{u}f(\vec{x}(u, v)) = (\nabla f) \cdot \pderiv{\vec{x}}{u} \\
  0 &= \pderiv{}{v}f(\vec{x}(u, v)) = (\nabla f) \cdot \pderiv{\vec{x}}{v}
\end{align*}
where the second equalities follow from the multivariate chain rule.

We know from IA Vector Calculus that $\pderiv{\vec{x}}{u}$ and $\pderiv{\vec{x}}{v}$ are linearly independent vectors tangent to $S$ and so since $\nabla f$ is orthogonal to both of them, $\nabla f$ is normal to $S$ at stationary points.
Therefore we have that $\vec{x}_0 \in S$ is a stationary point of $f$ \textbf{on $S$} if and only if $\nabla f(\vec{x}_0)$ is normal to the surface.

However, we often do not have an \textit{explicit} parameterisation of the surface and instead have a surface defined implicitly, for example $g(\vec{x}) = 0$ for some $g: \R^3 \to \R$.
Thus we are trying to extremise $f$ but now \textit{subject to the constraint} $g(\vec{x}) = 0$.
\begin{remark}[Recap]
  Recall from IA Vector Calculus that $\nabla g$ is normal to the surface $S$ at every point on the surface.
  This is because if we introduce a parameterisation $\vec{x}(u, v)$ which need not be specified, then by taking the partial derivatives of $g(\vec{x}(u, v)) = 0$ with respect to $u$ and $v$ we obtain:
  \[
    (\nabla g) \cdot \pderiv{\vec{x}}{u} = 0 \text{ and } (\nabla g) \cdot \pderiv{\vec{x}}{v} = 0
  \]
  at all points on the surface.
\end{remark}
Thus $\nabla f(\vec{x}_0)$ is normal to $S$ at $x_0$ if and only if it is parallel to $\nabla g(\vec{x}_0)$ and so $\vec{x}_0$ is a stationary point if and only if $\nabla f(\vec{x}_0) = \lambda \nabla g(\vec{x}_0)$ for some $\lambda \in \R$.
This means that extremising $f$ over $S$ amounts to solving this along with $g(\vec{x}_0) = 0$ to find $\vec{x}_0$ and $\lambda$.

Conveniently, we can obtain both of these equations by finding \textit{unconstrained} stationary points of $\varphi: \R^4 \to \R$ defined as:
\[
  \varphi(\vec{x}, \lambda) = f(\vec{x}) - \lambda g(\vec{x})
\]
This is because if we extremise with respect to $\lambda$:
\[
  \pderiv{\varphi}{\lambda} = 0 \iff g = 0
\]
and similarly if we extremise with respect to $\vec{x}$:
\[
  \nabla_{\vec{x}} \varphi = \vec{0} \iff \nabla_{\vec{x}} f - \lambda \nabla_{\vec{x}} g
\]
So finding stationary points of $\varphi$ over $\R^4$ is equivalent to extremising $f$ over the surface $S$.

We call $\lambda$ a \textit{Lagrange multiplier}.
This method can be generalised to $\R^{n}$ in a similar fashion.
We can also deal with multiple constraints by introducing more Lagrange multipliers.
For $m < n$ constraints $g_k(\vec{x}) = 0$ for $k \in \{1, \ldots, m\}$, we introduce $\vec{\lambda} = (\lambda_1, \ldots, \lambda_m) \in \R^{m}$ and then extremise:
\[
  \phi(\vec{x}, \vec{\lambda}) = f(\vec{x}) - \vec{\lambda} \cdot \vec{g}(\vec{x})
\]
with respect to $\vec{x}$ and $\vec{\lambda}$.
\section{Examples}
\subsection{Minimising Cuboid Surface Areas}
Consider a box with no lid with side lengths $x, y, z$.
We want to determine $(x, y, z)$ that minimises the surface area $A$ with fixed volume $V = L^3/2$.
Mathematically:
\begin{mini*}
  {\vec{x} \in \R^3}{A(\vec{x}) = 2z(x + y) + xy}{}{}
  \addConstraint{xyz}{= L^3/2}
\end{mini*}
\subsubsection{Using the Hessian}
We can solve the constraint directly to obtain that if $z = \frac{L^3}{2yz} \implies A(x, y) = \frac{L^3}{x} + \frac{L^3}{y} + xy$.
Stationary points of $A(x, y)$ are exactly those satisfying:
\begin{align*}
  0 &= \pderiv{A}{x} = - \frac{L^3}{x^2} + y \implies x^2 y = L^3 \\
  0 &= \pderiv{A}{y} = - \frac{L^3}{y^2} + x \implies y^2 x = L^3
\end{align*}
which has solution $x = y = L$.
Thus $z = L / 2$ and $A = 3L^2$.
To check this is a minimum point, by computing the hessian of $A(x, y)$, we see that is has positive eigenvalues and so it must be a minimum.
\subsubsection{Lagrange Multiplier Method}
Let:
\[
  \varphi(x, y, z, \lambda) = A(x, y, z) - \lambda\left(xyz - \frac{L^3}{2}\right)
\]
To extremise $\varphi$ :
\begin{align*}
  0 &= \pderiv{\varphi}{z} = 2(x + y) - \lambda xy \implies \lambda = \frac{2(x + y)}{xy} \\
  0 &= \pderiv{\varphi}{x} = 2x + y - \lambda yz \\
    &= 2z + y - \frac{2(x + y)z}{x} \\
    &= \frac{y}{x}(x - 2z) \implies x = 2z \text{ as $y = 0$ violates constraint} \\
  0 &= \pderiv{\varphi}{y} = 2z + x - \lambda xz \\
    &= \frac{x}{y}(y - 2z) \implies y = 2z \text{ as $x = 0$ violates constraint} \\
  0 &= \pderiv{\varphi}{\lambda} = xyz - \frac{L^3}{2} \implies L^3 = 2xyz
\end{align*}
Upon solving, we obtain $z = L/2$, $x = y = L$ and $\lambda = 4/L$, which is as before.

This disadvantage of this method is that there is no simple method for checking the nature of the stationary point.
In the above example, Hessian of $\phi$ is always a saddle point since $\pderiv[2]{\varphi}{\lambda} = 0$.
\subsubsection{Conclusion}
Although in this case it was easy to solve the constraint, this is often not the case and so the Lagrange multiplier method is usually preferable.
The Lagrange multiplier method is also often faster, even when the constraint can be solved.
\subsection{Minimising Quadratic Forms}
Suppose that we wish to find the minimum of the quadratic form $f(\vec{x}) = \vec{x}^{\trans} A \vec{x} = x_i A_{i j} x_j$ (noting the summation convention) on the surface $|\vec{x}|^2 = 1$ in $\R^{n}$.
Note that we can assume $A$ is symmetric here as any anti-symmetric component would not affect $f$.

We define:
\[
  \varphi(\vec{x}, \lambda) = x_i A_{i j} x_j - \lambda (|\vec{x}|^2 - 1)
\]
To find $\nabla_{\vec{x}} \varphi$, we compute its components:
\begin{align*}
  \pderiv{\varphi}{x_i} &= \pderiv{}{x_i}(x_k A_{k j} x_j - \lambda (x_j x_j - 1)) \\
                        &= A_{i j} x_j + A_{k i} x_k - 2\lambda x_i \\
                        &= 2(A_{i j} x_j - \lambda x_i) \\
                        &= 2[A \vec{x} - \lambda \vec{x}]_i \\
  \implies \nabla_{\vec{x}} \varphi &= 2(A\vec{x} - \lambda \vec{x})
\end{align*}
Finding the stationary points of $\varphi$, we require:
\begin{align*}
  \nabla_{\vec{x}} \varphi &= \vec{0} \iff A\vec{x} = \lambda \vec{x} \\
  \text{ and } \pderiv{\varphi}{\lambda} &= 0 \iff |\vec{x}|^2 = 1
\end{align*}
This means that the stationary points of $\varphi$ are exactly all the pairs $(\vec{x}, \lambda)$ of normalised eigenvectors and eigenvalues of $A$.

Evaluating $f$ at a stationary point $\vec{x}_0$ which is an eigenvector with eigenvalue $\lambda$, we see that:
\[
  f(\vec{x}_0) = x_i A_{i j} x_j = \lambda x_i x_i = \lambda
\]
Thus the eigenvalues of $A$ are the values of $f$ at stationary points and so the minimum of $f$ is the smallest eigenvalue of $A$.
\subsubsection{Alternative Approaches}
\begin{itemize}
  \item Use a rotation to diagonalise $A$ and allow us to rewrite $f$ as:
    \[
      f(\vec{x}) = \lambda_1 (x_1')^2 + \cdots + \lambda_n (x_n')^2
    \]
    noting that $|\vec{x}'| = |\vec{x}| = 1$.
    It is now easy to see that to minimise $f$, we want to direct all of $\vec{x}'$ along the direction of the smallest eigenvector, say $\lambda_k$.
    That is, we take $x'_i = 0$ for $i \neq k$ and $x'_k = 1$.
    We then find that the minimum of $f$ is $\lambda_k$ and can undo this rotation to recover the optimal $\vec{x}$.
  \item We can also consider the following ratio: \label{alternativeApproach}
    \[
      \Lambda(\vec{x}) = \frac{f(\vec{x})}{g(\vec{x})}
    \]
    where $g(\vec{x}) = |\vec{x}|^2$.
    We can rewrite this as:
    \[
      \Lambda(\vec{x}) = \hat{x}_i A_{i j} \hat{x}_j = f(\uvec{x}) \text{ where $\uvec{x} = \frac{\vec{x}}{|\vec{x}|}$}
    \]
    Therefore, the minimum of $f$ with respect to $\uvec{x}$ (i.e. constrained to the unit sphere) is equivalent to minimising $\Lambda(\vec{x})$ with respect to $\vec{x} \in \R^n$ (i.e. unconstrained).

    To find the extrema of $\Lambda$, since $\vec{x}$ is unconstrained here, we can just set $\nabla \Lambda = \vec{0}$:
    \[
      \vec{0} = \nabla \Lambda = \frac{1}{g}\left(\nabla f - \frac{f}{g} \nabla g\right) \iff 0 = \frac{2}{g}\left(A_{i j}x_j - \frac{f}{g}x_i\right)\ \forall i \iff A \vec{x} = \Lambda \vec{x}
    \]
    Thus the values of the extrema of $\Lambda$ are exactly the eigenvalues of $A$.
    Therefore by the above discussion, the extrema of $f$ over the unit sphere are also these eigenvalues.
\end{itemize}
\subsection{Probability Distributions}
What probability distribution $\{p_1, \ldots, p_n\}$ satisfying $\sum_{i = 1}^{n} p_i = 1$ and $p_i \geq 0\ \forall i$ maximises the \textit{information entropy} defined as:
\[
  S = - \sum_{i = 1}^{n} p_i \log_{2} p_i
\]
We define:
\[
  \varphi(p_1, \ldots, p_n, \lambda)= - \sum_{i = 1}^{n} p_i \log_{2} p_i - \lambda \left(\sum_{i = 1}^{n} p_i - 1\right)
\]
Recalling that $\log_{2}(p_i) = \frac{\log p_i}{\log 2}$, upon differentiating, we obtain:
\begin{align*}
  0 &= \pderiv{\varphi}{p_i} = - \log_{2}(p_i) - p_i \cdot \frac{1}{p_i \log(2)} - \lambda \iff p_i = 2^{(\frac{1}{\log 2} + \lambda)} \\
  0 &= \pderiv{\varphi}{\lambda} = \sum_{i = 1}^{n} p_i - 1 \iff \sum_{i = 1}^{n} p_i = 1
\end{align*}
This tells us that all $p_i$ must be equal and so $p_i = \frac{1}{n}\ \forall i \in \{1, \ldots, n\}$.
Thus the discrete uniform distribution maximises the information entropy.
Finally, substituting back in, we obtain that $S_{\text{max}} = \log_{2} n$.
\end{document}
